\section{Demo Scenario and Evidence}

The demo uses a small external Python task API repository with route handlers, a service layer, structured application errors, tests, and a short task specification. The specification acts as a task contract: it preserves response envelopes and handler signatures while leaving implementation strategy open. Before coding begins, the developer authors two freeform rules with \texttt{sc rules add}: one that compiles into a CLI-enforced constraint and one that compiles into prompt-level behavioral guidance. This setup makes the separation between deterministic governance, softer preferences, and task-specific specification visible before the main coding task starts.

The live workflow then proceeds in two short sessions. In session 1, Smart Coder receives the task and the spec, produces a spec-aware plan, and surfaces an API design tradeoff through a model-initiated check-in rather than guessing silently. The developer chooses an option and gives one explicit instruction about how subsequent low-risk work should proceed. In session 2, Smart Coder performs a related follow-up task on the same repository, now starting from loaded interaction history and reducing unnecessary interruption while preserving review at the API boundary. The demo closes with one observability view, showing that the run was fully traced and exported. This sequence highlights the system's intended contribution: adaptive oversight, not static prompting or unconstrained autonomy.

\begin{table}[t]
  \caption{Evidence to include in the final submission from one or more representative rehearsal runs.}
  \label{tab:evidence}
  \centering
  \begin{tabular}{p{0.44\columnwidth}p{0.42\columnwidth}}
    \toprule
    Evidence item & Current source \\
    \midrule
    Hard-rule compilation example & \textbf{TODO:} \texttt{sc rules add} output \\
    Guidance-rule compilation example & \textbf{TODO:} \texttt{sc rules add} output \\
    Representative task success & \textbf{TODO:} exported session bundle \\
    Session 1 model vs policy check-ins & \textbf{TODO:} \texttt{observe checkin-stats} \\
    Session 2 fewer interruptions than Session 1 & \textbf{TODO:} session trace comparison \\
    Session 2 starts from loaded history & \textbf{TODO:} run bootstrap output \\
    Verification outcomes & session bundle + trace CSV \\
    \bottomrule
  \end{tabular}
\end{table}

At minimum, the final submission should include one clean exported session bundle, one trace CSV, and a short video demonstrating the full two-session flow. Those artifacts support the three claims this paper needs: the system runs end to end on a realistic coding task, both model-side and policy-side oversight occur, and observed developer behavior changes later system behavior.
